%LTeX: language=fr-FR
\documentclass{article}
\usepackage{tasks}

\usepackage{amsmath,amssymb,amsfonts}
\usepackage{graphicx}
\usepackage{amsthm}
\usepackage[french]{babel}
\usepackage{enumitem}
\usepackage{currfile}

\usepackage{tikz}
\usetikzlibrary{calc,intersections}
\usepackage{qrcode}
\usepackage{mathrsfs}
\usepackage{enumitem}
\usepackage[french]{babel}
\usepackage{tcolorbox}
\usepackage{fancyhdr}
\usepackage{multicol}
\usepackage[headsep=3mm,footskip=5mm,includefoot,includehead]{geometry}
%showframe
\usepackage[lastpage,user]{zref}

\newtheorem*{theorem*}{Théorème}

\setlength{\parindent}{0cm}

\newenvironment{clist}{
  \begin{center}
  \begin{inparaenum}}{
    \end{inparaenum}
\end{center}}

\pagestyle{empty}
\graphicspath{ {../../Images/} }
\input{\currfiledir ../def.tex}
\def \Chapitre {Chapitre 2 -- Ensembles de nombres}
\def \ChapitreCourt {Ch. 2 -- Ensembles de nombres}

\geometry{left=1.5cm, right=1.5cm, top=1cm, bottom=1cm}
\setlength{\headheight}{14.0pt}

\theoremstyle{definition}
\newtheorem{exercice}{Exercice}

\begin{document}
\pagestyle{fancy}

\fancyhf{} % clear existing header/footer entries
% We don't need to specify the O coordinate
\fancyhead[L]{\large \textbf{\Niveau}}
\fancyhead[C]{\large \textbf{\Chapitre}}

\large

Au cycle 4, les élèves ont étudié les inégalités pour comparer des valeurs numériques. Ils ont résolu des équations et des
inéquations du premier degré. La notion d'intervalle, présentée comme ensemble de nombres vérifiant des inégalités, est
nouvelle. Elle prolonge le travail mené au cycle 4 sur les inéquations.

\begin{paragraph}{Pré-requis}
  \begin{itemize}
    \item Cycle 4 : inégalités
    \item inéquations du premier degré.
  \end{itemize}
\end{paragraph}

\begin{paragraph}{Contenu}
    \begin{itemize}
        \item Ensemble R des nombres réels, droite numérique.
        \item Intervalles de R. Représentation graphique, notations du type $[a, +\infty[$, $]-\infty, a]$, $[a, b]$
        \item Intersection et union de deux intervalles (vocabulaire ensembliste : $\cap$, $\cup$).
        \item Inégalités et opérations (ajout d'un nombre, multiplication par un nombre positif ou négatif).
        \item Ensemble des nombres décimaux. Encadrement décimal d'un nombre réel à $10^{-n}$ près.
        \item Ensemble Q des nombres rationnels. Nombres irrationnels ; exemples fournis par la géométrie, par exemple $\sqrt{2}$ et $\pi$.
    \end{itemize}
\end{paragraph}

\begin{paragraph}{Objectifs}
  \begin{itemize}
    \item Lire l'abscisse d'un nombre réel sur une droite graduée et placer un nombre réel d'abscisse donnée.
    \item Représenter un intervalle de la droite numérique. Déterminer si un nombre réel appartient à un intervalle donné.
    \item Donner un encadrement, d'amplitude donnée, d'un nombre réel par des décimaux.
    \item Dans le cadre de la résolution de problèmes, arrondir en donnant le nombre de chiffres significatifs adapté à la situation
étudiée.
  \end{itemize} 
\end{paragraph}

\begin{paragraph}{Démonstrations}
    \begin{itemize}
        \item $\dfrac{1}{3}$ n'est pas décimal : démontrée dans ce chapitre (raisonnement par l'absurde).
        \item $\sqrt{2}$ est irrationnel : énoncée dans ce chapitre, démontrée dans le chapitre d'arithmétique (fin d'année).
    \end{itemize}
\end{paragraph}

\begin{paragraph}{Approfondissements}
    \begin{itemize}
        \item Développement décimal illimité d'un nombre réel.
        \item Observation, sur des exemples, de la périodicité du développement décimal de nombres rationnels.
    \end{itemize}
\end{paragraph}

\begin{paragraph}{Exemple d'algorithme}
    \begin{itemize}
        \item Déterminer par balayage un encadrement de $\sqrt{2}$ d'amplitude inférieure ou égale à $10^{-n}$.
    \end{itemize}
\end{paragraph}

\end{document}
